\paragraph{When ICL emerges.} \citet{garg2022} trained transformers on inputs that varied equally in every direction, achieving near-least-squares performance on unseen linear functions. When test inputs instead varied much more in some directions than others, giving the model additional input--output examples of the same function eventually stopped reducing its prediction error on a new input. Least squares continued to improve using those same examples. In their test, the input covariance had eigenvalues proportional to $1/i^2$ in a random basis, and the model's error tracked least squares up to about ten examples and then plateaued, whereas noiseless least squares becomes exact once there are as many examples as dimensions. They also tested inputs from a low-dimensional subspace, noisy labels, rescaled prompts and queries in a different orthant from the examples, and found that bigger models reduced the error under these shifts considerably. They trained a GPT-2 \citep{radford2019gpt2} with Adam (learning rate $10^{-4}$, batch 64, 500,000 steps) and a curriculum over input dimension and prompt length. We take their model size, input dimension, prompt length and batch size, but put the same uneven input variances into training and vary the optimizer.

Which algorithm a trained model runs can be read off by comparing its predictions with reference algorithms \citep{akyurek2023}. Akyürek et al.\ constructed transformers that implement gradient descent and closed-form ridge regression, and found that in noiseless regression trained models moved from one step of gradient descent (one layer) through ridge regression to least squares (eight or more layers) as depth grew. \citet{bai2023} constructed transformers that implement least squares, ridge, Lasso, generalized linear models and gradient descent in context, and showed that one transformer can choose a different algorithm for each prompt. For large widths and depths the trained models approach Bayes-optimal estimators. A transformer can therefore represent many in-context algorithms, so which one a trained model runs depends on how it was trained. We use the same comparison, with principal component and ridge regression as references, because these are the two natural estimators when input directions differ in variance.

\citet{olsson2022} linked ICL in language models to induction heads, which form in an abrupt phase change visible as a bump in the loss: a short window early in training in which in-context learning improves sharply. \citet{raventos2023} found a task-diversity threshold: below it, a transformer acts as the Bayes estimator for its finite pretraining tasks; above it, it behaves like ridge regression and solves unseen tasks. Ridge corresponds to a Gaussian prior over all tasks, so above the threshold the model departs from the Bayes-optimal estimator for its own pretraining distribution; in their main setting the threshold lies between $2^{14}$ and $2^{15}$ tasks. More weight decay lowered the threshold, while more layers or a larger embedding raised it, so training choices move which estimator emerges. \citet{singh2023} showed that emergent ICL can be transient, giving way to in-weights learning with continued training in their setting, synthetic few-shot classification, even while the training loss keeps falling; more classes, Zipfian data and larger embeddings delayed the transience, and L2 regularization made ICL more persistent. Raventós et al.\ and Singh et al.\ trained with Adam, and Singh et al.\ list other optimizers as future work.

\paragraph{What attention computes, and how training gets there.} \citet{vonoswald2023} showed that a trained linear self-attention layer implements one step of gradient descent on the in-context regression loss: a single layer trained on isotropic regression converges to this construction, while deeper stacks learn a curvature-corrected variant. For inputs with covariance $\Sigma$, the best one-layer linear attention instead preconditions the step by approximately $\Sigma^{-1}$, with a finite-context correction $\mathrm{tr}\,\Sigma/n$ that acts as a regularizer \citep{ahn2023}. Gradient flow from a suitable initialization reaches this optimum, which predicts $\hat y = x_q^\top \Gamma^{-1} \frac1N\sum_i y_i x_i$ with $\Gamma = (1+\frac1N)\Lambda + \frac{\mathrm{tr}\,\Lambda}{N} I$ \citep{zhang2023trained}; the optimal gains $\rho_j^\star$ we use as a reference in linear attention (Section~\ref{sec:methods}) come from this solution. This end point shrinks the weak directions smoothly, as ridge regression does. The distribution of the inputs also shapes the learned algorithm more than the distribution of the targets given the inputs \citep{mahankali2023}. \citet{zhang2023trained} also tested shifts at evaluation time: the trained layer tolerates shifts of the task and, for long prompts, of the query, but not a rescaling of all test inputs, and GPT-2 models of Garg et al.'s size were more robust but still fell short of least squares at the largest rescaling. Our query-shift test is close to theirs, but we read it during training, at matched in-distribution error, and across optimizers. \citet{wang2025testtime} unified linear attention, state-space models and softmax attention as associative memories that solve a key-to-value regression in the forward pass. Three choices, the regression weights, the regressor class and the test-time optimizer, generate linear attention and its gated variants, state-space models such as Mamba-2, fast-weight programmers and softmax attention; plain linear attention is one gradient step that ignores the covariance between key dimensions.

These results describe the end point of training or what an architecture can represent; optimizers differ in the path. \citet{zhang2025training} analysed linear attention's training dynamics on this task: merged key--query weights learn in one abrupt drop, while separate ones learn one principal input direction at a time, so that a model stopped early approximately runs principal component regression (PCR) on a growing set of directions. Training goes from saddle to saddle, direction $m$ takes time proportional to $\lambda_m^{-2}$, each drop lowers the loss by about $\lambda_m$, and both parametrizations end at the same solution; one-layer softmax attention with separate rank-one keys and queries shows the same stages. This variance-ordered learning mirrors the exact solutions of \citet{saxe2014} for deep linear networks, in which input--output modes are learned in order of their strength. PCR uses the strongest directions fully and ignores the rest; ridge regression uses every direction, shrinking the weak ones more.

\paragraph{How optimizers shape learning.} SGD, Adam and Shampoo \citep{gupta2018shampoo} can be read as steepest descent under different norms \citep{bernstein2024old}. Without its moving averages Adam becomes sign descent, steepest descent under the $\ell_\infty$ norm of the parameters, and Shampoo becomes spectral descent, the update $-\eta\,UV^\top$ for a gradient $G = USV^\top$, which is steepest descent under the spectral norm. Muon \citep{jordan2024muon} applies this spectral step to the momentum of each hidden weight matrix, approximating $UV^\top$ with a few Newton--Schulz iterations and leaving embeddings, output head, gains and biases to AdamW. Adding decoupled weight decay and matching the update size across matrix shapes lets Muon scale to language-model pretraining with about twice AdamW's compute efficiency \citep{liu2025muon}. AdamW itself is Adam with weight decay decoupled from the gradient update \citep{loshchilov2019}; we set the decay to zero, so that no arm has an explicit regularizer, because weight decay alone moves which estimator emerges \citep{raventos2023}. \citet{xie2025adam} showed that Adam's advantage depends on the coordinate axes: rotating GPT-2's parameter-space objective slowed Adam sharply, below even a rotation-equivariant SGD variant, while SGD is unaffected by such rotations. On language data, gradient descent makes slow progress on rare classes where Adam and sign descent are much less affected \citep{kunstner2024heavy}; by analogy one might expect AdamW to have learned more of the weak input directions than SGD at the same error, but SGD has learned at least as much (Table~\ref{tab:h1}).

Several recent analyses ask which directions spectral optimizers learn first. \citet{ma2026} proved that simplified Muon trains linear transformers on in-context regression at a rate independent of the condition number, unlike gradient descent and SignGD. \citet{dragutinovic2026} showed that simplified Muon trades gradient descent's simplicity bias for speed in linear models and in linear attention on this task, where it learns every principal direction at the same rate (head $i$ finishes in time proportional to $\lambda_i^{-1/3}$ rather than one after another), with worse generalization in experiments where tasks shared structure, including a two-layer softmax transformer on a cycle task. By contrast, \citet{vasudeva2025} and \citet{wang2025muon} found that this more even learning helps spectral optimizers on rare classes: spectral descent learns all principal components of Gaussian-mixture data at equal rates, which improves class-balanced loss early in training, a gap that persists when gradient descent uses normalized step sizes and that depth amplifies. In language models, most of Muon's advantage over Adam sits in the attention value/output weights and the MLPs, and in a one-layer associative memory with imbalanced classes Muon learns all classes evenly while Adam's errors can be very uneven. The even learning is fragile: small-batch gradient noise pushes Muon's update back toward the gradient and restores strongest-directions-first learning in deep linear networks \citep{jacobs2026}, and depth amplifies the separation between modes. This matters here because we train with batches of 64. \citet{dong2026muonp} and \citet{shumaylov2026} studied updates that flatten the singular values only partly, judging them by loss and downstream accuracy, which Shumaylov et al.\ found barely depends on the update's spectrum: on GPT-2 the best Schatten exponents lie in a quasi-norm regime that no steepest-descent geometry represents, and replacing the singular values with random noise matches Muon, while zeroing the largest ones falls short. Dong and Sawin's update raises the singular values to a power $p$, interpolating between gradient descent ($p=1$) and Muon ($p=0$); it helps most in fine-tuning, while full flattening remains better for pretraining. We use the same knob to ask what the spectrum does to the learned algorithm rather than to the loss.

These optimizer studies motivate comparing learning under imbalanced training inputs, while Garg's result motivates testing robustness when the input distribution changes at evaluation.

\paragraph{Contributions.} The analyses above study simplified Muon (no momentum, exact orthogonalization) in linear models, linear attention or associative memories, mostly in the gradient-flow or full-batch limit, and judge optimizers mostly by speed or loss. We train SGD, AdamW and full Muon on in-context linear regression whose inputs vary much more in some directions than others, in GPT-2 models of Garg et al.'s size, a smaller GPT-2 and one-layer linear attention. We compare the optimizers at the same in-distribution error, so that speed is factored out, and measure how much of each input direction the model has learned and which in-context algorithm it runs. We shift test queries toward input directions that were rare or common in training to test when Muon helps out of distribution and when it hurts, and use a knob that flattens Muon's update spectrum by degrees to test why. Muon trains about three times faster and, at the same error, its in-context predictor is less PCR-like; out of distribution it helps on rare directions and hurts slightly on common ones; and flattening the update is what makes the difference.
